\begin{table}[H]
    \centering
    \caption{Muatan Faktor Terstandarisasi dan Komunalitas Subtes pada Model 2 (CFA)}
    \label{tab:factor_loadings}
    \begin{spacing}{1.0}
    \begin{tabular}{llcccc}
        \toprule
        Dimensi Laten & Subtes & $\lambda$ (Std.) & $SE$ & $z$ & $R^2$ ($h^2$) \\
        \midrule
        Literasi        & PPU          & 0,602$^{***}$ & 0,021 & 28,73 & 0,362 \\
        Literasi        & PBM          & 0,622$^{***}$ & 0,019 & 32,93 & 0,387 \\
        Literasi        & LBI Saintek  & 0,498$^{***}$ & 0,022 & 22,15 & 0,248 \\
        Literasi        & LBI Soshum   & 0,384$^{***}$ & 0,024 & 16,03 & 0,148 \\
        Literasi        & LBE          & 0,629$^{***}$ & 0,019 & 33,76 & 0,395 \\
        Numerasi        & PU           & 0,736$^{***}$ & 0,017 & 43,68 & 0,541 \\
        Numerasi        & PK           & 0,622$^{***}$ & 0,020 & 30,90 & 0,387 \\
        Numerasi        & PM           & 0,554$^{***}$ & 0,020 & 27,80 & 0,307 \\
        \bottomrule
    \end{tabular}
    \end{spacing}
    \vspace{2mm}
    \begin{flushleft}
    \footnotesize\textit{Catatan.} Bintang pada $\lambda$ menandakan signifikansi tiap muatan: $^{*}p < 0,05$; $^{**}p < 0,01$; $^{***}p < 0,001$. $SE$ = kesalahan standar; $z$ = nilai uji rasio kritis; $R^2$ ($h^2$) = proporsi varians yang dijelaskan oleh faktor laten terkait. Korelasi antarfaktor laten adalah $\phi = 0,813$$^{***}$.
    \end{flushleft}
\end{table}
